\section{An exact quantum comparison and the boundary measure}
\label{sec:quantummeasure}

The preceding question has a concrete partial answer. A special
\emph{unit-volume Liouville correlation} gives the Deligne--Mostow
K\"ahler potential exactly, at finite quantum parameter. Independently,
the boundary Reidemeister-torsion measure reduces to the Goldman measure,
whose rescaled flat limit is the Deligne--Mostow measure. Neither
statement identifies an arbitrary tetrahedral quantum amplitude with
the model integral~\eqref{eq:modelpartition}. The distinction between a
correlation, a K\"ahler potential, and a volume density is essential.

\subsection{A first-moment identity at the flat weights}

Fix the weights~\eqref{eq:flatweights}, all marked points finite in the
chosen coordinate, and a Liouville parameter $0<b<1$. Set
\begin{equation}
 \gamma=2b,\qquad Q=b+b^{-1},\qquad
 \beta_i=\frac{\mu_i}{b},\qquad
 s=\frac{\sum_i\beta_i-2Q}{\gamma}=-1.
 \label{eq:Liouvilleflatparameters}
\end{equation}
Here $\beta_i$ are the charge conventions of
Kupiainen--Rhodes--Vargas; in the physics convention
$V_{\alpha_i}=e^{2\alpha_i\phi}$, one has $\beta_i=2\alpha_i$.
Write $\Lambda_L>0$ for the cosmological constant, to distinguish it
from both the deficit weights and the geometric area.

The exact correlation formula of~\cite[eqs.~(2.11)--(2.15)]{KRV}
is
\begin{align}
 C_b(\boldsymbol\beta;\boldsymbol p)
   &=\frac1b\Lambda_L^{-s}\Gamma(s)
      P_b(\boldsymbol p)\,\mathbb E[Z_{\boldsymbol\beta}^{-s}],
       \label{eq:LiouvilleGMC}\\
 P_b(\boldsymbol p)
   &=\prod_{i<j}|p_i-p_j|^{-\beta_i\beta_j},\notag\\
 Z_{\boldsymbol\beta}
   &=\int_{\CC}\prod_i
       \left(\frac{|w|_+}{|w-p_i|}\right)^{\gamma\beta_i}
       M_\gamma(\dd^2w),\qquad |w|_+=\max(1,|w|),\notag
\end{align}
where $M_\gamma$ is the normalized Gaussian multiplicative-chaos
measure in their reference metric. The only probabilistic property
needed for the following calculation is its first moment,
\begin{equation}
 \mathbb E[M_\gamma(\dd^2w)]=|w|_+^{-4}\dd^2w.
 \label{eq:GMCfirstmoment}
\end{equation}
The standard unit-volume normalization removes the zero-mode
Gamma factor:
\begin{equation}
 C_b^{\mathrm{uv}}
       :=\frac{\Lambda_L^s}{\Gamma(s)}C_b
        =\frac1b P_b\,\mathbb E[Z_{\boldsymbol\beta}^{-s}].
 \label{eq:unitvolumeLiouville}
\end{equation}
The right side defines the continuation at $s=-1$.

\begin{proposition}[An exact Liouville--area identity]
For every $0<b<1$ and the fixed weights~\eqref{eq:flatweights},
\begin{equation}
 \boxed{\quad
 C_b^{\mathrm{uv}}(\boldsymbol p)
   =\frac1b\prod_{i<j}|p_i-p_j|^{-\mu_i\mu_j/b^2}
       \mathcal A_\mu(\boldsymbol p),\qquad
 \mathcal A_\mu=\int_{\CC}\prod_i|w-p_i|^{-2\mu_i}\dd^2w.
 \quad}
 \label{eq:exactLiouvilleDM}
\end{equation}
Consequently, on the configuration space away from collisions,
\begin{equation}
 \boxed{\quad
 -2i\,\partial\bar\partial\log C_b^{\mathrm{uv}}
       =\omega_{\mathrm{CH},\mu},\qquad
 \dd V_{\mathrm{CH},\mu}
       =\frac1{d!}\left(-2i\,\partial\bar\partial
                            \log C_b^{\mathrm{uv}}\right)^d,
 \quad d=n-3.}
 \label{eq:exactLiouvillecurvature}
\end{equation}
The form descends to $\M_{0,n}$ with its fixed weights and the
holomorphic-sectional-curvature-$-1$ convention.
\end{proposition}

\begin{proof}
At~\eqref{eq:Liouvilleflatparameters}, the exponent $-s$ is $1$.
Moreover, $\gamma\beta_i=2\mu_i$ and
$\gamma\sum_i\beta_i=4$. Applying~\eqref{eq:GMCfirstmoment}
therefore cancels the reference-metric factors exactly:
\[
 \mathbb E[Z_{\boldsymbol\beta}]
 =\int_{\CC}
    \prod_i\left(\frac{|w|_+}{|w-p_i|}\right)^{2\mu_i}
    |w|_+^{-4}\dd^2w
 =\mathcal A_\mu(\boldsymbol p).
\]
This integral is finite because each $\mu_i<1$ and, when all
marks are finite, the integrand is $O(|w|^{-4})$ at infinity.
It is exactly the area in~\eqref{eq:DMarea}.
Each $\log|p_i-p_j|$ is pluriharmonic off the diagonals.
Taking $-2i\partial\bar\partial\log$ consequently removes
$P_b/b$, leaving~\eqref{eq:CHpotential}. A change of the
M\"obius normalization changes the potential by a pluriharmonic
function, so its curvature descends.
\end{proof}

This is a deduction from the established correlation formula and
the area interpretation of the period pairing; it does not use
the main DOZZ structure-constant theorem. It is also an exact
finite-$b$ identity, rather than an interchange of a semiclassical
limit with two moduli derivatives. The raw expression
$C_b^{\mathrm{uv}}\dd^{2d}p$ is \emph{not} the moduli-space
volume density: the density is the top exterior power of its
logarithmic curvature.

\subsection{Analytic continuation and the equivalent residue}

The original unconstrained zero-mode integral is not convergent
at $s=-1$. The unit-volume function, however, is positive and
finite there. Indeed, the extended Seiberg bounds of
\cite[eq.~(2.14)]{KRV} require
\[
 -s<\min\left\{\frac4{\gamma^2},
       \min_i\frac2\gamma(Q-\beta_i)\right\},\qquad \beta_i<Q.
\]
At our parameters these inequalities reduce to
\begin{equation}
 1<b^{-2},\qquad
 1<1+\frac{1-\mu_i}{b^2},\qquad \beta_i<Q,
 \label{eq:flatSeibergbounds}
\end{equation}
and are strict. In particular, the positive moment has some
integrability beyond order one. Theorem~6.1 of~\cite{KRV}
gives holomorphic continuation in the charges of the unit-volume
$n$-point correlation to a complex neighborhood of this real
domain. Thus the pole in~\eqref{eq:LiouvilleGMC} is an actual
meromorphic pole, not just a formal screening calculation.

Since $\operatorname{Res}_{s=-1}\Gamma(s)=-1$,
\begin{equation}
 \operatorname*{Res}_{s=-1}C_b
      =-\frac{\Lambda_L}{b}P_b\mathcal A_\mu.
 \label{eq:onescreeningresidue}
\end{equation}
This residue means the coefficient of $(s+1)^{-1}$ with the
other charge coordinates fixed. With the transverse physics
coordinate $t=\sum_i\alpha_i-Q+b=b(s+1)$, it becomes
$\operatorname{Res}_{t=0}C_b=-\Lambda_LP_b\mathcal A_\mu$.
Consequently the logarithmic curvature of the absolute residue
also equals~\eqref{eq:exactLiouvillecurvature}.

In Coulomb-gas terminology this is the one-screening resonance
$\sum_i\alpha_i=Q-b$; see also
\cite[Section~3.4]{TeschnerRevisited}. Its chiral contour integrals are
\begin{equation}
 \int_\gamma\prod_i(w-p_i)^{-\mu_i}\dd w,
 \label{eq:screeningDMperiods}
\end{equation}
precisely the Deligne--Mostow periods~\eqref{eq:DMperiods}.
Their Hermitian pairing is the area. Section~\ref{sec:screenedfusion}
gives an independent finite period calculation of the four-point
residue and fusion matrix, and states the additional continuation
and tail-control input for the spectral-contour derivation. A pointwise
identification with an unnormalized modular-double $6j$ function
would still require its normalization and contour to be specified.
Section~\ref{sec:normalizedsixj} fixes these normalizations
for the induced finite operator and examines the geometric
meaning of that specialization.

\subsection{Four points: the exact quantum potential and its measure}

Put the marks at $(0,1,z,\infty)$, with the standard primary-field
normalization at infinity. Then~\eqref{eq:exactLiouvilleDM} reads
\begin{align}
 C_b^{\mathrm{uv}}(z)
   &=\frac1b |z|^{-\mu_1\mu_3/b^2}
              |1-z|^{-\mu_2\mu_3/b^2}\mathcal A_\mu(z),
              \label{eq:fourpointLiouville}\\
 \mathcal A_\mu(z)
   &=\int_{\CC}|w|^{-2\mu_1}|w-1|^{-2\mu_2}
                         |w-z|^{-2\mu_3}\dd^2w.\notag
\end{align}
The condition at infinity is now
$\mu_1+\mu_2+\mu_3=2-\mu_4>1$. The ordinary field
normalization at infinity cancels the corresponding area scaling;
there is no additional moduli-dependent factor.

For equal weights $\mu_i=1/2$, the area is
\begin{equation}
 \mathcal A(z)=\int_{\CC}\frac{\dd^2w}{|w(w-1)(w-z)|}.
 \label{eq:fourpointareaintegral}
\end{equation}
Let $\Omega_1,\Omega_2$ be oriented lattice periods of
$\dd w/y$ on the elliptic curve~\eqref{eq:ellipticDM}, and put
$\tau=\Omega_2/\Omega_1=u+iv\in\HH$. The pillowcase has
half the area of its double cover, so
\[
 \mathcal A(z)=\frac12\operatorname{Im}
       (\overline{\Omega_1}\Omega_2)
             =\frac12|\Omega_1|^2v.
\]
Locally $\Omega_1$ is nonzero and holomorphic. Hence
\begin{equation}
 \boxed{\quad
 -2i\partial\bar\partial\log C_b^{\mathrm{uv}}
    =\frac{\dd u\wedge\dd v}{v^2}
    =\frac{|\tau'(z)|^2}{(\operatorname{Im}\tau(z))^2}
                         \dd(\operatorname{Re}z)\wedge
                         \dd(\operatorname{Im}z).
 \quad}
 \label{eq:fourpointquantumDMdensity}
\end{equation}
The elliptic-integral expression for $\tau$ is
\eqref{eq:ellipticperiodratio}. Integrating this form over the
labeled moduli space gives $2\pi$, as in
Section~\ref{sec:thurstonDM}. This is the flat equal-deficit
Deligne--Mostow area; the cusped hyperbolic WP area is the
different quantity $V_{0,4}(0)=2\pi^2$.

\subsection{An actual boundary torsion measure}

There is also a measure comparison on the representation side.
For a good, boundary-regular $\mathrm{SL}_2(\CC)$ representation
of an $n$-holed sphere, Heusener--Porti prove
\begin{equation}
 \Omega_S=\pm\frac{\omega_G^d}{d!}
                  \wedge\nu_1\wedge\cdots\wedge\nu_n,
                  \qquad d=n-3,
 \label{eq:torsionGoldman}
\end{equation}
where $\Omega_S$ is the holomorphic Reidemeister-torsion volume
form, $\omega_G$ is the Goldman form on the fixed-conjugacy
leaf, and the $\nu_i$ are the peripheral volume forms
\cite[Theorem~3]{HeusenerPorti}. With invariant pairing
$B(U,V)=-\tr(UV)$ their peripheral formula is
$\nu_i=\pm\sqrt2\,\dd\tr\rho(\partial_i)$
\cite[Proposition~2]{HeusenerPorti}.
The pairing must be the same in all three kinds of form.
For $B\mapsto\lambda^2B$, with $\lambda\in\CC^*$,
the homogeneity stated after that theorem is, in our notation,
\[
 \omega_G\mapsto\lambda^2\omega_G,\qquad
 \nu_i\mapsto\lambda\nu_i,\qquad
 \Omega_S\mapsto\lambda^{2d+n}\Omega_S.
\]
Thus the displayed $\sqrt2$ belongs specifically to $B=-\tr$;
it is not retained unchanged when the Goldman form is recalibrated.
If the scale of $B$ is negative on the real form, choose its
complex square root consistently in the holomorphic forms,
then take the positive density on the real leaf.

One must disintegrate by the peripheral forms \emph{before}
fixing the conjugacy classes. Simply pulling the unrestricted
top-degree form back to a fixed-trace leaf gives zero. After
this disintegration, and on the relevant real component, the
torsion density is the Goldman Liouville density, with the
chosen normalization of $B$.

Mondello identifies the cone WP and Goldman Poisson structures
up to a specified constant for openings in $(0,2\pi)$
\cite[Theorems~1.1--1.2]{MondelloCones}. We calibrate $B$ and
the torsion and peripheral forms together so that the symplectic
form on our real leaf is $\omega_{\WP}$, with the positive
trace-coordinate density fixed in Appendix~\ref{app:fricke}.
The holonomy statement is local in general; it is globally
injective on marked cone metrics when every opening is at
most $\pi$. Combining~\eqref{eq:torsionGoldman}
with~\eqref{eq:localWPCHlimit} gives, locally on a controlled
family,
\begin{equation}
 \boxed{\quad
 \left(\frac2A\right)^d
        \dd V_{\mathrm{torsion,relative}}
    \longrightarrow \dd V_{\mathrm{CH},\mu}.
 \quad}
 \label{eq:torsionDMlimit}
\end{equation}
Apply the torsion theorem for $A>0$ and then take the rescaled
limit. At the Euclidean endpoint the affine holonomies become
reducible, and their ordinary characters discard translational
shape data. The limiting twisted periods or a blowup retaining
their scale are indispensable. In particular, an unqualified
ordinary Goldman construction for the Euclidean isometry group
would miss the needed nondegenerate pairing.

\subsection{The four-point character surface, including a correction}

The preceding measure limit can be calculated directly. Give
the four cone points the common opening
$\theta=\pi-A/4$, so their total hyperbolic area is $A$.
Choose equal peripheral traces
$a=-2\cos(\theta/2)=-2\sin(A/8)$. The Fricke relation
of Appendix~\ref{app:fricke} becomes
\begin{equation}
 \Phi=x^2+y^2+z^2+xyz-2a^2(x+y+z)+4a^2+a^4-4=0.
 \label{eq:equalFricke}
\end{equation}
Blow up the limiting character by
\begin{equation}
 x=-2-AX,\qquad y=-2-AY,\qquad z=-2-AZ.
 \label{eq:Frickeblowup}
\end{equation}
Direct substitution gives
\begin{align}
 \frac\Phi{A^2}
  ={}&X^2+Y^2+Z^2-2XY-2XZ-2YZ+1\notag\\
     &+A\left(-XYZ+\frac{X+Y+Z}{8}\right)
       -\frac{A^2}{768}+O(A^3).
 \label{eq:Frickeexpanded}
\end{align}
On a chart with $u>0$, $v>0$, parametrize the limiting quadric by
\begin{equation}
 X=\frac1{2v},\qquad
 Y=\frac{u^2+v^2}{2v},\qquad
 Z_0=\frac{(u-1)^2+v^2}{2v}.
 \label{eq:Frickeperiodchart}
\end{equation}
These are squared side lengths of a Euclidean triangle of area
$1/4$; the quadric equation is Heron's identity. Since
\[
 X+Y-Z_0=\frac uv,\qquad
 \dd X\wedge\dd Y=\frac{u}{2v^3}\dd u\wedge\dd v,
\]
the calibrated density~\eqref{eq:frickedensity} yields
$A^{-1}\omega_{\WP}\to\frac12\dd u\wedge\dd v/v^2$,
with the positive orientation on this chart.

In fact, the implicit-function expansion is
\[
 Z=Z_0+AZ_1+O(A^2),\qquad
 Z_1=\frac{(1-u)[v^2-u(1-u)]}{16v^2}.
\]
Using
$\Phi_z=2A(X+Y-Z)+A^2(XY-1/8)+O(A^4)$ gives
\begin{equation}
 \boxed{\quad
 \frac{\omega_{\WP}}A
  =\frac1{2v^2}\left[
       1-\frac{A(1+u^2+v^2)}{16v}+O(A^2)\right]
                    \dd u\wedge\dd v.
 \quad}
 \label{eq:Frickefirstcorrection}
\end{equation}
This is an explicit first correction in a fixed \emph{blown-up
character chart}. Here $\tau=u+iv$ is the limiting flat period
coordinate, not the exact period ratio or conformal cross-ratio
of the finite-$A$ surface. The error is local on compact subsets
of the chart. In particular, the displayed correction is not
termwise integrable at $v=\infty$; it cannot be integrated
naively to recover the global polynomial of
\eqref{eq:exactWPvolumeflat}.

\subsection{What a normalized quantum kernel would have to produce}

The phase-space measure of a unitary change of polarization is
more constrained than a bare scalar $6j$ prefactor. In one
degree of freedom, let $S(q,Q)$ generate a canonical relation,
$p=\partial_qS$, away from its caustic. A branchwise normalized
semiclassical kernel has the leading form
\begin{equation}
 K_\hbar(q,Q)\sim
  \frac{|\partial_q\partial_QS|^{1/2}}{\sqrt{2\pi\hbar}}
       \exp\left(\frac{iS}{\hbar}+\frac{i\pi m}{4}\right).
 \label{eq:normalizedWKBkernel}
\end{equation}
The amplitude follows by applying unitarity to the leading
phase difference $(q-q')p$ and changing variables from $Q$ to
$p$. Thus
\begin{equation}
 2\pi\hbar|K_\hbar|^2\dd q\dd Q
     \sim |\partial_q\partial_QS|\dd q\dd Q
     =|\dd q\wedge\dd p|.
 \label{eq:kernelLiouvillemeasure}
\end{equation}
This is a branchwise assertion. Several stationary branches
produce interference; a pointwise square of an oscillatory
cosine need not converge without averaging.

Teschner--Vartanov's Fenchel--Nielsen transition kernel contains
the factors $\sqrt{M(\alpha_s)M(\alpha_t)}$ multiplying the
$b$-$6j$ symbol~\cite[eq.~(5.14)]{TeschnerVartanov}.
Consequently a proposed measure comparison must include the
spectral measure, label Jacobians, and polarization. If a
normalized kernel quantizes our cone WP leaf, its classical
measure is constrained by~\eqref{eq:kernelLiouvillemeasure},
and its rescaled flat measure by~\eqref{eq:torsionDMlimit}.
This conditional argument does not by itself establish the
elliptic continuation of that kernel.

\subsection{The auxiliary tetrahedron and the order of limits}

There is a concrete reason not to identify the auxiliary
tetrahedron in a fusion kernel with the convex realization of
our cone sphere without further work. In the conventions of
\cite[Sections~3--5]{TeschnerVartanov}, the parameters satisfy
\begin{equation}
 \alpha_k=\frac Q2+\frac{i\ell_k}{4\pi b},\qquad
 A_k=e^{-2\pi ib\alpha_k+\pi i}
           =e^{\ell_k/2-\pi ib^2}.
 \label{eq:b6jparameterdictionary}
\end{equation}
Here $A_k$ is the paper's exponentiated tetrahedral parameter,
not our geometric area $A$. Continuing a boundary length by
$\ell_i=i\theta_i$ gives an auxiliary dihedral parameter
$\eta_i=\theta_i/2-\pi b^2$, modulo the chosen branch.
Thus at the equal-angle Euclidean endpoint, and as $b\to0$,
the four external auxiliary dihedrals tend to $\pi/2$.
The internal $s,t$ edges are opposite. After reordering faces,
the limiting angle Gram matrix would be
\begin{equation}
 G=\begin{pmatrix}
 1&-\cos\eta_s&0&0\\
 -\cos\eta_s&1&0&0\\
 0&0&1&-\cos\eta_t\\
 0&0&-\cos\eta_t&1
 \end{pmatrix},\qquad
 \det G=\sin^2\eta_s\sin^2\eta_t.
 \label{eq:auxiliaryGramobstruction}
\end{equation}
For interior angles $\eta_s,\eta_t\in(0,\pi)$ this is
nonsingular. A nondegenerate Euclidean tetrahedron has a
rank-three angle Gram matrix. Hence this real slice cannot
converge to the generic Euclidean tetrahedron realizing our
flat cone metric. An additional degeneration, a different
real slice, or a different geometric identification is needed.
This is an obstruction to the proposed identification, not
to the existence of Liouville theory with cone charges.

There is also a finite-$b$ shift in the charge convention.
At a fixed classical opening $\theta_i=2\pi(1-\mu_i)$, the
exact dictionary~\eqref{eq:b6jparameterdictionary} gives
\[
 \alpha_i=\frac{\mu_i}{2b}+\frac b2,
\]
whereas~\eqref{eq:Liouvilleflatparameters} uses
$\alpha_i=\mu_i/(2b)$. The latter corresponds in that
dictionary to the shifted opening $\theta_i+2\pi b^2$.
For unshifted flat openings the zero-mode parameter would be
$s=n/2-1$, not $-1$. Thus the exact correlation identity is
a specified charge specialization; it is not an identification
with the finite-$b$ kernel at unshifted classical cone angles.
Indeed, treating the shifted openings as hyperbolic cone data
would give the negative Gauss--Bonnet area $-2\pi n b^2$.

There is a second, independent distinction. Along a heavy-charge
path with deficit weights $a_i(A)$, set
$\alpha_i(A)=a_i(A)/(2b)$ and
$A=2\pi(\sum_i a_i(A)-2)$. Then exactly
\begin{equation}
 \boxed{\quad s=\frac{A}{4\pi b^2}-1.\quad}
 \label{eq:areaLiouvillecharge}
\end{equation}
The finite-$b$ flat endpoint of this pure-heavy charge path is
$s=-1$, the resonance used
above. The original zero-mode convergence region $s>0$
requires $A>4\pi b^2$. In contrast, under the proposed joint
scaling $A\asymp\hbar^{2/3}$ with $\hbar\asymp b^2$, one
has $A/b^2\to\infty$ and hence $s\to+\infty$.
The exact residue identity therefore establishes a quantum
origin for the Deligne--Mostow curvature without proving the
uniform quantum degeneration proposed in
Section~\ref{sec:quantumdegeneration}.

